\subsection{Higher vanishing moments}
We use the arbitrary-rank higher-indicator theorem and its twist-sector
corollary from \cite[Theorem~\ref{B-thm:higher-indicator} and
Corollary~\ref{B-cor:higher-sector}]{Finite}. For
$W_\zeta(V)=\FPdim(I(V)_\zeta)$ they state that
\[
 (r+1)W_\zeta(V)\leq D\FPdim(V)
 \quad\text{if }\nu_1(V)=\cdots=\nu_r(V)=0.
\]
Here $\mathcal C$ is pseudo-unitary with its canonical spherical structure.
The proof uses the nonnegative polynomial
$|1+z+\cdots+z^r|^2$ and the moment identity
\begin{equation}\label{eq:indicator-moments}
 \sum_Z[I(V):Z]\FPdim(Z)\theta_Z^n=D\nu_n(V).
\end{equation}
In particular, non-self-dual simples have vanishing first and second
indicators and satisfy
\begin{equation}\label{eq:nonselfdual-third}
 3\FPdim(I_{\rm low}(X))\leq D\FPdim(X).
\end{equation}

\begin{remark}
The proof also gives an explicit bound when the indicators do not vanish:
\[
 (r+1)^2W_\zeta(V)\leq
 (r+1)D\operatorname{FPdim}(V)
 +2D\operatorname{Re}\sum_{j=1}^r(r+1-j)\zeta^{-j}\nu_j(V).
\]
Only positivity of the canonical dimensions and the indicator formula are used;
no $*$-structure is assumed. Positivity must not be dropped in a general spherical
category. The polynomial is the unnormalised Fej\'er kernel.
\end{remark}

\begin{corollary}[Aggregating nontrivial simples]\label{cor:aggregate}
Let $V=\bigoplus_{i>0}a_iX_i$, with $a_i\in\mathbb Z_{\geq0}$, and let
$L_V$ be a known contribution to its twist-one mass. Then
\[
 2L_V\leq D\FPdim(V),\qquad
 9L_V\leq D\bigl(3\FPdim(V)+2\sum_{i>0}a_i\bigr).
\]
\end{corollary}
\begin{proof}
The first indicator vanishes on every summand, while
$\nu_2(V)=\sum a_i\nu_2(X_i)\leq\sum a_i$. Apply the degree-one and
degree-two Fej\'er polynomials. Equivalently, add the individual bounds.
\end{proof}
